% !TeX root = ../papers/nc-p3.tex
\section{四元群の一文字次数が欠ける片側対称窓}
\label{sec:klein-companion-windows}
$K=(\mathbb Z/2)^2$を文字次数群とし、片側対称な二写像
$\beta:E^*\otimes F^*\to H,\ \gamma:E^*\otimes G^*\to H^*$が斉次とする。
三つの点空間は各次数一次元の正則表現とする。
テンソル直線の指標は必要なら$G$の文字次数を平行移動して吸収する。
隣接階数は6である。以下で被覆を与えるのは、$H$で少なくとも一つの文字次数が欠ける場合である。
単純な固定点とは、四元群で固定される共通点でその辺の零点が横断的なものをいう。

\begin{theorem}\label{thm:klein-missing-character}
上の欠ける次数を持つ全入力で、標準AS正則性を有限に判定できる。
座標固定点の一つが単純な標準AS窓は、両対称チャートまたは前節の三軸族に属する。
特に欠ける次数を初めから仮定せずとも、単純な固定点のあるこの斉次類の
非両対称AS窓は全て、三係数の非零な正則三軸族で尽くされる。
全ての座標固定点が非横断的な場合には、後述する二平面ずつのPBW表示または
両対称チャートにより判定できる。
\end{theorem}

\paragraph{文字ブロック。}
定理\ref{thm:single-opposite-evaluation-ranks}の座標文字での四評価下界と
両関係空間に純テンソルがないことは必要である。以後この有限条件を満たすとする。
欠ける次数を平行移動して零とし、
$b_{ab}=\beta(e_a,f_b), c_{ab}=\gamma(e_a,g_b)$と書く。
$b_{aa}=c_{aa}=0$、$a\ne b$では全て非零である。
両対称性を課すのは偶側だけで、式は
$\langle b_{ab},c_{cd}\rangle=\langle b_{cb},c_{ad}\rangle$である。
従って$c_{ba}$は$b_{ab}$を消し、その非零な次数空間は二次元以上である。
総次元6より三つの非零次数は各二次元である。
一色の四ベクトル$b_{ab},b_{ba},b_{cd},b_{dc}$のうち三つが比例すれば、
その直線を消す共変ベクトルが一成分だけの純関係を与えるので、三つは比例しない。
$c_{ab}$の直線は$b_{ba}$の零化直線で尽くされる。

欠ける次数を仮定せずに単純な固定点から始めた場合もこの表示に入る。
偶座標点の唯一の奇核は固定された座標点である。二つの軸評価の像は、
それぞれ同じ一つの次数を除く三つの次数に直線を持つ。横断性から各次数で
この二直線は独立である。従って残る一つの次数は零である。

\paragraph{単純な頂点。}
$0$を単純な座標点とし、$U=\beta(e_0,F^*)$と$V=\beta(E^*,f_0)$を取ると
$H=U\oplus V$である。$\gamma(e_0,G^*)=V^\perp$、
$\gamma(E^*,g_0)=U^\perp$である。三つの軸評価で$u\in V,y\in U,z\in U^*$を
補座標に取ると、最後の軸は$V\to V^*$の可逆対称写像となる
（$f_0,g_0$での偶対称性）。異なる文字次数が直交するのでその対角を1に揃えられる。
両写像は
\[
 \beta=(r y+A_u y,\;s u+B_y u),\qquad
 \gamma=(r z+A_u^{\mathsf t}z,\;t u+C_z u)
\]
となる。転置の式は$e_0$を入れて得られ、$B_y,C_z$は$g_0,f_0$を入れて対称となる。
斉次性から、巡回順$i,j,k$で
$A_i e_j=\alpha_i e_k, A_i e_k=\delta_i e_j, A_i e_i=0$であり、
$B_i=b_i(E_{jk}+E_{kj}), C_i=c_i(E_{jk}+E_{kj})$である
（以下の$b_i$は一つの係数で、先の$b_{ab}$とは区別する）。
残る偶対称式は
\begin{equation}\label{eq:klein-flat-coefficients}
 b_i c_j=-\delta_i\delta_j,\qquad b_j c_i=-\alpha_i\alpha_j
 \qquad((i,j,k)\text{ 巡回})
\end{equation}
で尽くされる。これは$\beta(e_i,f_i)$と$\gamma(e_j,g_j)$の対合を、
偶文字を交換した値と比較し、奇文字も交換した比較を加えたものである。
その他の比較は文字次数で消えるか同じ式になる。

全係数が非零なら、三式を掛けて$\prod\alpha_i/\prod\delta_i=\epsilon=\pm1$を得る。
$U$の三座標を対角に変えると$\delta_i=\epsilon\alpha_i$にできる。
実際、各辺で対角倍率の平方の比を指定する三条件の積が1であり、代数閉体上で平方根が取れる。
変更後の六式から$b_i=d\alpha_i,c_i=-d^{-1}\alpha_i$となる。
$y,z$の相反する共通倍率で$d$を吸収する。
$\epsilon=-1$はまさに前節の三軸族である。$\epsilon=1$なら$A_i$は全て対称で、
$\gamma((r,u),(-s,y))=\operatorname{diag}(I,-I)\beta((r,u),(s,y))$となるため
両対称判定に帰着する。

係数が零になる場合も残さない。例えば$b_i=0$なら奇座標$f_i$への評価の
上側だけで階数3が必要なので$\delta_j\alpha_k\ne0$である。
式\eqref{eq:klein-flat-coefficients}から$\delta_i=\alpha_i=0$となる。
今度は$e_i$での両評価が下側だけで階数3となるため
$b_jb_kc_jc_k\ne0$、従って$c_i=0$である。
$c_i=0$から始めても転置で同じである。$b,c$が全て非零なら六式から他も全て非零になる。
残る二つの$A$は文字の三頂点上で二辺だけを支えるため、両方を対角変更で同時に対称にできる。
六式で$c_j/b_j=c_k/b_k$となるので、適切な奇単位の倍率と
$H$の下側計量の倍率を選ぶと再び両対称表示を得る。
以上で単純な頂点に関する被覆が従う。
\begin{remark}
全係数非零の$\epsilon=1$側も常に標準ASである。$L=\operatorname{diag}(\alpha_i)$として
上側の積は$D(Lu,y)$である。両核は
$u_i^2-u_j^2+(\alpha_j^2-\alpha_i^2)r^2$の張る同じ鉛筆である
（反対切断の平方以外の三ブロックは可逆二行交代行列、平方ブロックは階数2）。
三平方が相異なるときは対角完全交叉が滑らかである。対角鉛筆の特異な各元は
単独の座標頂点を頂点に持ち、その点はもう一式で消えないためである。
純関係もない。上側・下側の非零な二係数ベクトルをともに使えば、
零の$00$成分を挟む二行二列で階数2以上である。片側だけなら非零の
$D$の極行列が非零対角倍率を除いて含まれ、その非零成分の二行二列小行列式がある。
定理\ref{thm:smooth-apolar-complete}を適用する。
二平方が等しければ符号変更で$\alpha_i=\alpha_j$と取れる
（同時に三文字の対角符号と軸の倍率を変えればよい）。
因子$e_i+e_j,e_i-e_j$には逆順の同じ奇因子が対応する。
各三平面像は$e_i\pm e_j$の上側と下側および
$(e_k,e_k)$を張る（上下の非零な共通倍率を許す）。
よって交わりは直線であり、純関係がないので共通根基の二文字関係は可逆である。
定理\ref{thm:radical-factor-amalgam}からこの場合も正則となる。
\end{remark}

\paragraph{横断する頂点がない場合。}
四頂点の辺$ab$を$b_{ab},b_{ba}$が比例するとき選ぶ。頂点での二軸の交わりの次元は
そこに出入りする選択辺の本数である。従って各頂点に選択辺があれば、
互いに素な選択二辺があるか、選択辺が三辺の星をなすかのいずれかである。

\begin{lemma}[相補的な二つの選択辺]\label{lem:klein-matching-pbw}
相補的な二辺が選ばれた場合、座標評価下界と純関係の不在の下で窓は標準Koszul AS正則である。
\end{lemma}
\begin{proof}
二辺を$01,23$に取り、文字を$\{0,1\}$と$\{2,3\}$の二組に分ける。
この色の関係は各組に一本ずつの二項交換関係となる。
実際、二組がそれぞれ同一直線に写り、二直線は独立で全係数が非零である。
相手窓でも零化直線の表示により同じである。
残る四色成分の関係（他の二色で各二式）の、後の組を左、先の組を右に持つ
四単項式への射影は可逆である。例えば$b_{20},b_{31}$の独立性は
$\langle b_{20},c_{13}\rangle\ne0$と同値である。偶位置交換で
これは$\langle b_{10},c_{23}\rangle$となり、独立な二辺の直線の一つと
他方の零化共変ベクトルの対合だから非零である。$b_{30},b_{21}$、逆向きの二組、
$c$側も同じ理由による。従って四つの混合逆順を正順へ全て交換できる。
組内も二項式で整列する順序では、二次関係は全ての六逆順を先頭語とする。
次数3の正規候補語は20個である。一方テンソルの一文字切断は全階数である
（異なる次数の四評価が非零で、六次元への相手も全射）。
その四つの縮約は三次交叉に入るので三次商の次元は20以上となる。
従って三次で全ての重なりが解消され、全次数で二次Gr\"obner基底である。
このKoszul性と標準Hilbert級数、Gramテンソルの切断の完全対合から標準AS性を得る。
\end{proof}

最後に星なら、三色とも一方の辺が選択されている。相異なる奇添字$b,d$について
$m_{bd}=w_{abcd}$と置く。ここで$a,c$は残る二文字で、偶交換により順序によらない。
色$i$の二辺を$ab,cd$とすると、$b_{ab},b_{ba}$から$c_{cd},c_{dc}$への対合の
二行二列行列は、どちらかの辺が選択されていれば階数1である。
四成分は全て非零である。実際、選択辺の直線に残り一文字の値が入ることは
三つの比例による純関係障害で排除される。逆向きの二行二列行列も同様である。
これらは$m_{ij}$のうち相異なる四添字の二行二列式の全六本にあたる。
従って$m_{bd}=u_b v_d$ ($b\ne d$)と分離できる。
実際、行・列を頂点とする二部グラフ$K_{4,4}$から対角を除くと立方体の辺となり、
六つの小行列式は各面一周の倍率が1という条件である。六面の境界は全ての閉路を
整数係数で生成するので、任意の道で行・列の倍率を決められる。
奇文字$d$の一側を$u_d/v_d$倍すれば$m_{bd}$は$b,d$で対称である。
このほかの非零成分は偶・奇の各平方の対合だけなので、両対称表示を得る。
これで定理\ref{thm:klein-missing-character}の残りも証明された。

\begin{proposition}[三軸族の外の反対階数10]\label{ex:klein-matched-rank-ten}
交代行列$m$の上三角成分を辞書順に$(2,1,1,3,3,5)$とし、
偶平方・奇平方の対合を
\[
 T=(w_{abab})_{a,b}=
 \begin{pmatrix}0&0&1&1\\0&0&-35/2&-65/2\\
 2/15&21&0&0\\-2/15&-39&0&0\end{pmatrix}
\]
とする。四添字が相異なるときは$w_{abcd}=m_{bd}$、ほかは零とする。
三つの隣接色ブロックは各々階数2である。実際、辺を$ab=0h,cd$とし
$cd$を昇順に取ると、順序$ab,ba,cd,dc$でブロックは
$\left(\begin{smallmatrix}U&P\\-P^{\mathsf t}&V\end{smallmatrix}\right)$であり、
$P=m_{(b,a),(d,c)}$である。$h=1$なら$U=V=0$、$P$は全成分非零の階数1である。
$h=2,3$では$U$が可逆で$V=-P^{\mathsf t}U^{-1}P$となる。
各方向の関係行列で六逆順への射影が可逆であり、組内二関係は二項式である。
補題\ref{lem:klein-matching-pbw}の整列証明から全次数で標準ASとなる。
反対切断では平方ブロックの行列式が36で、残る六つの交代成分が独立だから階数10である。
各組の座標直線上では相手の同じ直線との間に非退化な二項式が一本だけ残る。
従って偶点スキームはねじれの位置にある二直線を含む。
三軸族の反対階数10では、全ての曲線成分が一つの平面内にある
（節\ref{sec:cross-product-parameters}）。この例は印を忘れても三軸族に入らない。
\end{proposition}

\begin{proposition}[次数次元が$(3,1,1,1)$の場合]\label{prop:klein-three-singletons}
四次数が全て現れ、その一つが三次元ならば、この斉次類の標準AS窓は
全対称表示を持つ。従ってこの型の任意の入力には全対称な完全交叉の判定が使える。
\end{proposition}
\begin{proof}
三次元の次数を零に移す。相異なる$a,b$について$b_{ab}=0$と仮定する。
残りの二つの$c$各々に$d=a+b+c$を取ると、偶対称式は
$0=\langle b_{cb},c_{ad}\rangle$となる。第$b$列の評価階数が3以上なので
両方の$b_{cb}$は非零である。この次数空間は一次元なので両方の$c_{ad}$が零となり、
$\gamma$の第$a$行の階数に反する。転置でも同じであり、全ての非対角成分は非零である。
各一文字空間の基底と双対基底を選ぶ。固定した$d$に対し残る$a,b,c$を巡回させて
$c_{ad}/c_{cd}=b_{ab}/b_{cb}$を掛けると、
$b_{ab}b_{bc}b_{ca}=b_{ba}b_{cb}b_{ac}$となる。
従って辺の比$b_{ab}/b_{ba}=u_a/u_b$と書ける。奇軸$f_b$を$u_b$倍すると
$\beta$は二文字に対して対称となる。零次数の対角成分はこの交換で動かない。
偶対称性に加えて隣接交換も得たので、系\ref{cor:two-positive-as-recognition}と
系\ref{cor:partial-symmetric-as-complete}を適用できる。
\end{proof}
